% Candidate D.76. Notes only pending full independent audits.
\subsection{The joint contribution of all three peripheral axes}

\begin{statement}{D.76}{Proved}{Three explicit compact factors for every abelian pair}
Let \(s=(a,b)\) be any abelian pair of the tripod over \(\mathbb Q\),
with arbitrary ramification, and let \(\mathcal K_s\) be the compact
term of D.63. For a completed Kummer class \(c\) define
\[
 \delta(c)=
 \begin{cases}
 0,&c=0,\\
 N(c)-r(c),&c\ne0,\quad N(c)<\infty,\\
 \infty,&c\ne0,\quad N(c)=\infty,
 \end{cases}                                               \tag{D.76.1}
\]
where \(r(c)\) and \(N(c)\) are the global and local-at-five
divisibility exponents of D.75. Put
\[
                     d(s)=\delta(a)+\delta(b)+\delta(a-b).
                                                               \tag{D.76.2}
\]
There is a continuous split surjection of compact modules
\[
 \mathcal K_s\twoheadrightarrow
 \begin{cases}
       \mathbb Z_5^{\,d(s)},&d(s)<\infty,\\
       \displaystyle\prod_{\mathbb N}\mathbb Z_5,&d(s)=\infty.
 \end{cases}                                               \tag{D.76.3}
\]
Thus the three contributions add: they are not merely three
possibly overlapping lower bounds.
The divisibility exponents are computed by the convergent
logarithm formula (D.75.3).

In particular \(s=(\kappa(101),0)\) has a split
\(\mathbb Z_5^2\) factor, and a cusp parameter of global
depth \(r\) and finite local depth \(N\) has a split
\(\mathbb Z_5^{2(N-r)}\) factor. The explicit logarithm
cancellation of D.75 still gives a countable product.
This theorem applies even when the metabelian section fibre
above \(s\) is empty. It asserts neither existence of such
a section nor a calculation of the complementary compact factor.
\end{statement}

\begin{proof}
\paragraph{All three are actual arithmetic quotients.}
Write \(\Lambda=\mathbb Z_5[[u,v]]\), \(t_1=1+u\),
\(t_2=1+v\), and
\[
 J_1=(v),\qquad J_2=(u),\qquad J_3=(t_1t_2-1).
 \]
Each is stable under the actual Alexander action.
The first quotient of \(A_s=\Lambda C\) is exactly
D.75's \(M_a\), for arbitrary \(b\): the factor \(t_2^{b(g)}\)
becomes one. The second quotient is \(M_b\).

For the third put \(t_1=t\), \(t_2=t^{-1}\).
The actual reflection identity of D.39 gives
\[
 \mathcal B_g(t,t^{-1})
             =\frac{\chi\,t^{(\chi-1)/2}}{S_\chi(t)},
 \qquad
 S_\chi(t^{-1})=t^{1-\chi}S_\chi(t).
 \]
Consequently the full Fox multiplier on that quotient is
\[
              \chi S_\chi(t)\,
                   t^{\,a(g)-b(g)+(1-\chi(g))/2}.       \tag{D.76.4}
\]
Since \(1/2\in\mathbb Z_5\), the extra exponent is an
integral coboundary and is removed by an integral translation.
Thus the third quotient is precisely \(M_{a-b}\), up to
the indicated actual gauge.

Let \(D_i\) be the Cartier duals of these three quotients,
viewed as submodules of \(D_s=A_s^D\).
The one-axis argument of D.75 therefore computes
\[
 Q_i:=D_i^{G_{\mathbb Q_5}}/D_i^{G_{\mathbb Q}}
 \simeq
 \begin{cases}
  (\mathbb Q_5/\mathbb Z_5)^{\delta(c_i)},&
                         \delta(c_i)<\infty,\\
  \displaystyle\bigoplus_{\mathbb N}\mathbb Q_5/\mathbb Z_5,&
                         \delta(c_i)=\infty,
 \end{cases}
 \quad(c_1,c_2,c_3)=(a,b,a-b).                           \tag{D.76.5}
\]
When \(c_i=0\), its global and local affine images are
the same unit group after one integral gauge, so \(Q_i=0\).
It remains to show that these quotients inject jointly.

\paragraph{The precise relation among the three axis duals.}
Consider the compact map
\[
 \Lambda\longrightarrow
       \Lambda/J_1\oplus\Lambda/J_2\oplus\Lambda/J_3.
                                                               \tag{D.76.6}
\]
Use \(u\) as coordinate on the first and third factors
and \(v\) on the second. A triple of power series
\((F(u),G(v),H(u))\) lies in its image exactly when
\[
 F(0)=G(0)=H(0),\qquad
                      H_1=F_1-G_1.                    \tag{D.76.7}
\]
Here the subscript denotes the linear coefficient.
Necessity follows from \(v=-u/(1+u)\) on the third axis.
For sufficiency, start with \(F(u)+G(v)-F(0)\).
Its discrepancy on the diagonal has zero constant and
linear coefficients. On that diagonal
\[
                         uv=-\frac{u^2}{1+u}.
 \]
Adding \(uv\,L(u)\) therefore corrects any such discrepancy
with an integral power series \(L(u)\).
This proves (D.76.7), and the cokernel of (D.76.6) is
exactly a free \(\mathbb Z_5\)-module of rank three.
All its defining functionals use only constant and linear
coefficients on each axis.

Duality now shows that the kernel of
\[
                       D_1\oplus D_2\oplus D_3
                             \longrightarrow D_s             \tag{D.76.8}
\]
has, in each component, a functional involving only the
constant and linear coefficients of the corresponding
one-variable module. Under the measure model of D.75,
such a functional is represented by a polynomial of degree
at most two in \(z\), modulo constant functions. Indeed
\[
       \frac{t^z-1}{t-1}
                =z+\binom z2(t-1)+\text{higher powers}.
                                                               \tag{D.76.9}
\]
Integral translation gauges replace a function by
\(f(z+d)-f(d)\), and so preserve this degree bound.
This remains true for the half-cyclotomic gauge in (D.76.4).
No independent-summands assertion has been assumed.

\paragraph{A polynomial defect cannot survive the tame unit.}
Suppose \(x_i\in D_i^{G_{\mathbb Q_5}}\) and that their
sum in \(D_s\) is globally invariant.
For each nonzero \(c_i\), make the compatible global and
local normalizations of D.75. Its global affine image is
still \(5^{r(c_i)}\mathbb Z_5\rtimes\mathbb Z_5^\times\),
and its local image has pure unit multipliers.
Represent \(x_i\) by a continuous function \(f_i\) modulo
constants, normalized by \(f_i(0)=0\).
Local invariance then gives
\[
                           f_i(2z)=f_i(z).             \tag{D.76.10}
\]
For \(\tau\in\ker\chi\), global invariance of the sum
puts the three component defects in the kernel of (D.76.8).
Writing \(h=c_i(\tau)\), and replacing \(\tau\) by its
inverse if necessary to match the dual-action convention,
we obtain
\[
 f_i(z+h)-f_i(z)=P_h(z)
                       =A_hz^2+B_hz+C_h,               \tag{D.76.11}
\]
with coefficients in \(\mathbb Q_5/\mathbb Z_5\).
A constant ambiguity is absorbed in \(C_h\).
The coefficients are unique: evaluation at \(0,1,2\)
has Vandermonde determinant two, a five-adic unit.

Equation (D.76.10) and the telescoping identity give
\[
              P_{2h}(2z)=P_h(z),\qquad
              P_{2h}(z)=P_h(z+h)+P_h(z).                \tag{D.76.12}
\]
The quadratic coefficient is therefore simultaneously
\(A_{2h}=2A_h\) and \(4A_{2h}=A_h\).
Thus \(7A_h=0\), whence \(A_h=0\).
The linear coefficients then give \(B_{2h}=2B_h\) and
\(2B_{2h}=B_h\), so \(3B_h=0\) and \(B_h=0\).
Finally the constants give \(C_{2h}=2C_h=C_h\),
so \(C_h=0\). Both seven and three are units on the
five-primary torsion target. Thus every defect (D.76.11)
is zero.

By the exact primitive-restriction argument of D.75,
the values \(h=c_i(\ker\chi)\) range over all of
\(5^{r(c_i)}\mathbb Z_5\). Hence \(f_i\) is fixed by
every global translation. It is also fixed by every pure
unit, from its local invariance and evaluation at zero.
The exact global affine image now makes \(x_i\) globally
invariant. If \(c_i=0\), this conclusion already holds
because the global and local affine images agree.
We have proved the genuine injection
\[
               Q_1\oplus Q_2\oplus Q_3
                  \lhook\joinrel\longrightarrow
                 D_s^{G_{\mathbb Q_5}}/D_s^{G_{\mathbb Q}}.
                                                               \tag{D.76.13}
\]

\paragraph{The compact direct factor.}
Compose the dual of (D.76.13) with D.63's continuous
surjection from \(\mathcal K_s\) onto the dual of its
right-hand side. Formula (D.76.5) gives exactly (D.76.3).
As in D.75, the direct sum on the left of (D.76.13) is
divisible, hence injective as a discrete \(\mathbb Z_5\)-module.
Its inclusion in \(\mathcal K_s^D\) has a retraction;
duality makes the compact surjection split continuously.
This determines the displayed factor only.
For a cusp pair \((a,0)\), the two equal nonzero Kummer
parameters \(a\) and \(a-0\) occupy different axes, and
the proved joint injection counts both of them.
This gives the final examples and explains why their ranks
double without an unsupported independence assumption.
\end{proof}
